\section*{Chapter \ref{ch:g11:public-choice-first-look} --- Public Choice: A First Look}
% figdata: figdata/grade-11/09-public-choice-first-look/out/ (sugar, sugar_key); Brenland numbers invented.

\begin{solution}{exo:g11:public-choice-first-look:1}
$600\,000\,000/3\,000\,000 = 200$ dalens a year.
\end{solution}

\begin{solution}{exo:g11:public-choice-first-look:2}
$300\,000\,000/1500 = 200\,000$ dalens each, $1000$ times the cost to a household.
\end{solution}

\begin{solution}{exo:g11:public-choice-first-look:3}
$123\,900/6.89 \approx 18\,000$ times. The benefit to producers includes the processors' gain, so dividing it by the number of farms gives more than a grower received.
\end{solution}

\begin{solution}{exo:g11:public-choice-first-look:4}
Steel tariff: concentrated benefits to domestic steel makers and their workers, dispersed costs to the buyers of steel and of goods made with it. Hairdressing licence:
benefits to licensed hairdressers, who face fewer rivals; costs to all their customers. Aircraft subsidy: benefits to the few makers; costs to all taxpayers.
\end{solution}

\begin{solution}{exo:g11:public-choice-first-look:5}
$50\,000/10\,000\,000 = 0.005$ dalens. Ten hours of study cost 200 dalens: not worth it on that ground alone.
\end{solution}

\begin{solution}{exo:g11:public-choice-first-look:6}
Each spends $V/4 = 200$ thousand dalens; together 400 thousand, half the licence's value.
\end{solution}

\begin{solution}{exo:g11:public-choice-first-look:7}
For $n = 2$: $V \times 1/4$, as found. The total is $n \times V(n-1)/n^2 = V(n-1)/n$: $3V/4 = 600$ thousand for $n = 4$ and $9V/10 = 720$ thousand for $n = 10$. As $n$
grows, the total approaches the whole value of the licence.
\end{solution}

\begin{solution}{exo:g11:public-choice-first-look:8}
$1000 - 200p = 100$ gives $p = 4.5$; consumption 450, production 350. Growers gain $\tfrac12(100 + 350) \times 2.5 = 562.5$ million dalens. The \omterm{def:g10:what-prices-do:dwl}{deadweight loss} is
$\tfrac12 \times 250 \times 2.5 + \tfrac12 \times 250 \times 2.5 = 625$ million.
\end{solution}

\begin{solution}{exo:g11:public-choice-first-look:9}
No. Learning could change her vote, but her vote changes the outcome with probability one in a million, so the expected value of learning is at most about $30 \times
10^{-6}$ dalens a year against a cost of 200. She would learn if the knowledge were useful for other reasons (she could buy less sugar, say), if learning were cheap, or if
she joined others to act together, which raises the effect of her effort.
\end{solution}

\begin{solution}{exo:g11:public-choice-first-look:10}
$a = \tfrac12(100 + 300) \times 2 = 400$; $b = \tfrac12 \times 200 \times 2 = 200$; $c = 200 \times 2 = 400$; $d = \tfrac12 \times 200 \times 2 = 200$ (million dalens).
The buyers' loss is $\tfrac12(700 + 500) \times 2 = 1200 = 400 + 200 + 400 + 200$.
\end{solution}

\begin{solution}{exo:g11:public-choice-first-look:11}
The assumption that voters are rationally ignorant of small costs. Once voters know, a representative who backs the programme risks votes, and the model predicts it
becomes harder to keep or is cut.
\end{solution}

\begin{solution}{exo:g11:public-choice-first-look:12}
A policy with \omterm{def:g11:public-choice-first-look:concentrated-benefits}{concentrated benefits and dispersed costs} that is ended while its beneficiaries still oppose it. In trucking the carriers holding route rights and the
unionised drivers gained, and both strongly opposed the end of the regulation; yet it ended in 1980. The model leaves out the shippers, firms that bought transport and bore
a concentrated part of its cost; the published evidence on what regulation cost, which made the issue salient; and a period of high inflation in which cutting prices
mattered to voters.
\end{solution}

\begin{solution}{exo:g11:public-choice-first-look:13}
The job pays $150 - 60 = 90$ thousand a year more, 450 thousand over five years. The ruling raises the chance of getting it by $0.4 - 0.1 = 0.3$, so it is worth
$0.3 \times 450 = 135$ thousand dalens to him in expectation. A rule barring officials from jobs in the industry for some years after leaving, or decisions taken by panels
that rotate, would weaken the incentive.
\end{solution}

\begin{solution}{exo:g11:public-choice-first-look:14}
No tax: a loss of 100 million. Tax: a loss of $100 + 30 = 130$ million. Here the market failure left alone costs less than the correction. The comparison depends on how
far lobbying can push the tax from the right level and on how much the contest costs; with a tax set within, say, a quarter of the right level, the tax would win.
\end{solution}

\begin{solution}{exo:g11:public-choice-first-look:15}
Alone, X gives net gains $+50, -20, -20$: two of three vote against it, and it fails; Y fails the same way. If 1 and 2 vote for both, both pass: each of them nets
$80 + 10 - 60 = 30$, and legislator 3 nets $20 - 60 = -40$. Here the outcome is efficient: each project brings 100 in benefits for 90 in costs. With benefits of 70, 10 and
0, the same deal would pass two projects bringing 80 for 90: vote trading can pass inefficient projects too.
\end{solution}

\begin{solution}{pb:g11:public-choice-first-look:1}
\textbf{1.} Consumption $900 - 200 = 700$, production $200 - 100 = 100$, imports 600 thousand tonnes.

\textbf{2.} Supply starts at $p = 1$: $\tfrac12 \times 100 \times (2 - 1) = 50$ million dalens.

\textbf{3.} Demand reaches zero at $p = 9$: $\tfrac12 \times 700 \times (9 - 2) = 2450$ million dalens.

\textbf{4.} $(900 - 100p) - (100p - 100) = 200$ gives $p = 4$.

\textbf{5.} Consumption 500, production 300 thousand tonnes.

\textbf{6.} $a = \tfrac12(100 + 300) \times 2 = 400$ million dalens.

\textbf{7.} $c = 200 \times (4 - 2) = 400$ million dalens.

\textbf{8.} $\tfrac12(700 + 500) \times 2 = 1200$ million dalens.

\textbf{9.} $b = \tfrac12 \times 200 \times 2 = 200$ and $d = \tfrac12 \times 200 \times 2 = 200$: 400 million. $400 + 400 + 400 = 1200$.

\textbf{10.} $b$, $d$ and $c$: 800 million dalens.

\textbf{11.} $400\,000\,000/200 = 2$ million dalens per grower; $1\,200\,000\,000/3\,000\,000 = 400$ dalens per household.

\textbf{12.} 5000 times.

\textbf{13.} $40/400 = 10\,\%$.

\textbf{14.} No. Her letter barely changes the chance that the quota ends, and if it ended she would gain at most 400 dalens a year: the expected gain is far below 20. Each
household gains if all write, but each would rather the others did: ending the quota is a \omterm{def:g11:public-goods-and-commons:private-good}{public good} for the buyers, and they free ride.

\textbf{15.} The prize is the quota rent, 400 million; two equal contestants spend half of it, 200 million.

\textbf{16.} $400 + 40 + 200 = 640$ million dalens a year.

\textbf{17.} $(400 - 40)/1200 = 30\,\%$.

\textbf{18.} The licence holders would pay their rent to the government instead of spending it on the contest: the 200 million of contest spending disappears, and the 400
million of rent becomes revenue. The full cost falls to $400 + 40 = 440$ million.

\textbf{19.} For example: printing the cost of the quota on each bag of sugar, or publishing each year's cost per household.

\textbf{20.} Each grower gains 2 million dalens a year while each household loses 400, five thousand times less; the growers spend 10\,\% of their gain on lobbying to keep
the quota.
\end{solution}
